% !TeX spellcheck = es_VE
\documentclass[12pt,letterpaper]{article}
\usepackage[utf8]{inputenc}
\usepackage{amsmath}
\usepackage{amsfonts}
\usepackage{amssymb}
\usepackage{graphicx}
\usepackage{url}
\usepackage{tikz}
\usepackage[shortlabels]{enumitem} % Se mantiene aquí con sus opciones
\usepackage{fancyhdr}
\usepackage[spanish]{babel}
%\usepackage[acronym]{glossaries}
%\makeglossaries

% agregar circuitos electronicos
\usepackage[siunitx, american]{circuitikz}

% comas sencillas en listing
\usepackage{textcomp}
\usepackage{upquote}
\usepackage{listings}
\lstset{upquote=true}

% dimensiones de pagina
\usepackage[left=1.5cm,top=1.5cm,right=1.5cm,bottom=2.5cm]{geometry}

% hiperenlaces
\usepackage{hyperref} % Para incluir hipervínculos
\hypersetup{urlcolor=blue,colorlinks=false,pdfborder={0 0 0}} % eliminar marco rojo en enlaces (citas, indices, pies de pagina, etc)

% permitir en la numeracion de los items agregar el numero de seccion
% (Se eliminó la línea duplicada \usepackage{enumitem} que causaba el error)
\setlist[enumerate]{label=\thesection.\arabic*. }   % numerar solo con seccion
%\setlist[enumerate]{label=\thesubsection.\arabic*. } % incluir numeracion de subsection

% otras configuraciones
\usepackage{parskip}
\setlength{\parskip}{1em} % establece el espacio entre parrafos
\setlength{\parindent}{0cm} % elimina la sangria inicial en parrafo
\usepackage{array} % ajustes adicionales para tablas
% \usepackage{pdfpages} %incluir archivos pdf en documento

% personalizacion de etiquetas
\renewcommand{\figurename}{Figura}
\renewcommand{\contentsname}{Contenido}
\renewcommand{\listfigurename}{Índice de Figuras}
\renewcommand{\listtablename}{Índice de Tablas}
\renewcommand{\lstlistlistingname}{Índice de Códigos}
\renewcommand{\tablename}{Tabla}
\renewcommand{\lstlistingname}{Código}

% pie de pagina
\renewcommand{\footrulewidth}{0.4pt}% grosor de linea del pie de pag
%\lfoot{\thepage} % texto izquierda
%\rfoot{\thepage} % texto derecha
\cfoot{\thepage} % texto al centro

%% definiciones para insertar codigo de programacion al doc
\usepackage{color}
\definecolor{dkgreen}{rgb}{0,0.6,0}
\definecolor{gray}{rgb}{0.5,0.5,0.5}
\definecolor{mauve}{rgb}{0.58,0,0.82}
\lstset{frame=tb,
    language=Python,
    aboveskip=3mm,
    belowskip=3mm,
    showstringspaces=false,
    columns=flexible,
    basicstyle={\small\ttfamily},
    %numbers=left,
    numberstyle=\tiny\color{gray},
    keywordstyle=\color{blue},
    commentstyle=\color{dkgreen},
    stringstyle=\color{mauve},
    breaklines=true,
    breakatwhitespace=true,
    tabsize=3
} % fin de codigo de programacion

%%%%%%%%%%%%%% MACRO PARA CONTROL DE PUNTAJE %%%%%%%%%%%%%%%%%%
\usepackage{xifthen}
\newcounter{puntos}
\setcounter{puntos}{0} % Inicializa el contador a cero

\newcommand{\sumarPunto}[1]{%
    \addtocounter{puntos}{#1}%
    \ifthenelse{#1 > 1}{(\textbf{#1 puntos})}{(\textbf{#1 punto})} 
}
%%%%%%%%%%%%%%%%%%%%%%%%%%%%%%%%%%%%%%%%%%%%%%%%%%%%%%%%%%